\documentclass[11pt,a4paper]{article}

\usepackage[T1]{fontenc}
\usepackage[utf8]{inputenc}
\usepackage[a4paper,margin=24mm]{geometry}
\usepackage{amsmath,amssymb}
\usepackage{url}
\Urlmuskip=0mu plus 1mu

\setlength{\parskip}{0.25em}
\setlength{\parindent}{1em}
\emergencystretch=3em
\widowpenalty=10000
\clubpenalty=10000
\displaywidowpenalty=10000

\newcommand{\effect}[1]{\mbox{\texttt{#1}}}
\newcommand{\meet}{\mathbin{\sqcap}}
\newcommand{\seq}{\mathbin{\cdot}}
\newcommand{\sem}[1]{[\![#1]\!]}
\title{Configurable Static Type Checking for Forth}
\author{Jaanus Pöial\\Tallinn University of Technology}
\date{}

\begin{document}
\pagenumbering{gobble}
\maketitle

\begin{abstract}
Forth stack comments document the cells consumed and produced by a word.
Ordinary comments are not checked and do not record how stack words
move or copy individual cells.  This paper presents a configurable source-level
static type checker implemented in gforth.  It describes each word by a typed 
input/output stack contract, called a stack effect.  A validation profile defines the type 
hierarchy and primitive word effects (specifications). Symbolic indices record 
the flow introduced by stack shuffling words such as
\texttt{DUP}, \texttt{SWAP}, and \texttt{OVER}.

The analysis uses three algebraic operations.  Composition handles linear word
sequences.  Meet operation unifies alternative branches.  A finite
one-pass-versus-two-pass comparison handles repetition. Forth examples provided in the paper
illustrate these operations.  They cover type refinement, cell identity,
branch merging, reusable inferred effects, and local
loop stability. The prototype derives every reported effect without executing concrete
values. Its conclusions depend on the user-defined type hirarchy, primitive
specifications and source text.
\end{abstract}

\noindent\textbf{Keywords:} Forth; static type checking; stack effects;
symbolic execution.

\section{Introduction}

A Forth word is a named operation.  Here, a cell is one single-cell stack
item.  The data stack is the ordered sequence of those items.  A word normally
receives operands from that stack and returns its results to the same stack.
The conventional notation
\begin{verbatim}
+   ( n1 n2 -- sum )
\end{verbatim}
places the top of stack on the right and says that \texttt{+} consumes two
numbers and produces one.  Such comments are valuable interfaces, but a Forth
system usually consumes them as source text rather than treating them as
proof obligations.  A comment can therefore be missing, stale, or semantically
stronger than the checked definition.

Consider the familiar definition
\begin{verbatim}
: SQUARED  ( n -- n^2 )  DUP * ;
\end{verbatim}
The checker can derive \texttt{( n -- n )} from the declarations of
\texttt{DUP} and \texttt{*}.  It can verify the number and identity of the
cells used by the phrase.  It does not prove the mathematical assertion that
the result is the square of the input.

\begin{samepage}
The checker reads three text inputs:
\begin{enumerate}
  \item a file that defines types, aliases, and ordering relations (type hierarchy);
  \item a file that defines word effects and source-reading rules;
  \item a Forth source file whose phrases are checked and whose colon
        definitions are inferred.
\end{enumerate}
A \emph{primitive word} is a word whose effect comes directly from the second
file.  A \emph{colon definition} begins with \texttt{:} and ends with
\texttt{;}.  The word that it defines receives an inferred effect from the
source between those delimiters.  A top-level phrase is checked and traced but
does not add a dictionary entry.
\end{samepage}
The prototype consumes the three files and maps each known word to one stack
effect.  At \texttt{;}, the inferred body effect is added to this dictionary.
The prototype analyzes symbolic stack cells; it does not execute the Forth
definition. The inferred definitions are stored in a log file.

The design builds on earlier work on stack effects and symbolic type variables
\cite{poial1994,poial1992,poial2002,poial2006}.
The following topics are discussed:
\begin{itemize}
  \item how a later word can constrain an earlier input;
  \item how two branch effects are combined;
  \item why an effect displays only the top cells that it touches;
  \item why the loop check compares one pass with two passes.
\end{itemize}
Sections~\ref{sec:model}--\ref{sec:control} define the analysis used here.
Sections~\ref{sec:sequence-examples}--\ref{sec:loop-examples} derive all twelve
test results and the four helper effects.

\section{Typed Symbolic Stack Effects}
\label{sec:model}

\subsection{Types and compatibility}

Let $T$ be a finite set of abstract cell types with a precision order
$\preceq$.  The relation $t\preceq u$ means that $t$ is more specific than
$u$.  A value of type $t$ may therefore be used where $u$ is required.  This
order is the profile's subtype relation.  The Forth-2012-like profile used for the examples contains, among
other edges, the transitive chains
\[
 \texttt{a-addr}\preceq\texttt{c-addr}\preceq
 \texttt{addr}\preceq\texttt{u}\preceq\texttt{n}\preceq\texttt{x},
\]
\[
 \texttt{char}\preceq\texttt{+n}\preceq
 \texttt{u}\preceq\texttt{n}\preceq\texttt{x},
 \qquad \texttt{flag}\preceq\texttt{x}.
\]
The flag type is deliberately separate from the numeric chain in this profile.
Another profile may classify flags as numbers.  The hierarchy is therefore a
policy for one target vocabulary.  It is not a universal claim about Forth
cells.  Names such as \texttt{x}, \texttt{X}, and \texttt{cell} may be aliases
for one type.

Two types are \emph{comparable} when one is below the other in $\preceq$.
Comparable types are compatible.  Matching them retains the more specific
type:
\[
 \operatorname{match}(t,u)=
 \begin{cases}
 t,&t\preceq u,\\
 u,&u\preceq t,\\
 \text{undefined},&\text{otherwise}.
 \end{cases}
\]
The prototype does not search for a third common subtype when $t$ and $u$ are
incomparable.  An undefined match is a type clash.

\subsection{Symbols, effects, and the implicit frame}

A symbolic cell is a type with an optional positive identity index, written
$t[i]$.  Equal indices denote the same abstract origin.  Different indices do
not claim that the concrete values differ.  Every unindexed occurrence starts
as an independent symbol.  An unindexed output is therefore a fresh abstract
result.  Indices exist only in the checker.  They are not addresses or runtime
tags.

\begin{samepage}
The declarations of three stack words illustrate the notation:
\begin{verbatim}
DUP   ( x[1] -- x[1] x[1] )
SWAP  ( x[2] x[1] -- x[1] x[2] )
OVER  ( x[2] x[1] -- x[2] x[1] x[2] )
\end{verbatim}
Before each word occurrence is analyzed, its local indices receive fresh
numbers.  Two primitive calls therefore cannot share an index by accident.
After analysis, \emph{normalization} renumbers the surviving indices from one
upward.  Renumbering itself changes only their printed names.  The
prototype also compacts identities: a derived index that records only an
implicit one-to-one correlation may be hidden, while an explicitly declared
or multiply occurring correlation is retained.  Once hidden, that pairwise
correlation is unavailable to later composition.  This deliberate precision
loss keeps stored effects small but means that inferred effects need not be
principal.
\end{samepage}

An effect $e=(L--R)$ contains two finite symbolic stacks, with their tops on
the right.  A stack \emph{suffix} is the group of top cells displayed by the
effect.  Deeper cells form an implicit \emph{frame}.  For every frame $F$, the
contract acts as
\[
                  F\,L \longmapsto F\,R.
\]
For example, \effect{( a-addr -- x )} for \texttt{@} does not require the
entire data stack to contain one cell.  It requires only that the top cell be
an aligned address; any deeper $F$ passes through unchanged.  The identity
effect is \effect{( -- )}.  Write $\epsilon$ for the empty stack.  A clash is
a separate failure result and must not be confused with the identity effect.
The set of symbolic effects is the checker's \emph{abstract domain}.  An
effect operation is \emph{partial} when some arguments produce no effect.  In
those cases, this checker reports a clash.

\section{Main Effect Operations}
\label{sec:operations}

\subsection{Sequential composition}

Write typed symbolic cells by juxtaposition, with the top on the right, and
let $e$ and $f$ denote possibly empty sequences of such symbols.  The two
symbolic equations from the presentation capture the basic cases:
\[
 (a -- b)\seq(e\,b -- d)=(e\,a -- d),
 \qquad
 (a -- f\,c)\seq(c -- d)=(a -- f\,d).
\]
In the first equation, the second effect consumes $b$ but also needs the deeper
sequence $e$; that sequence therefore becomes part of the composite input.  In
the second, the first effect leaves $f$ below $c$; consuming $c$ leaves $f$
below the new result $d$.  In case of type mismatch the result is zero (type clash).

More generally, $p=(L_p -- R_p)$ followed by $q=(L_q -- R_q)$ is written
$p\seq q$.  Boundary cells are matched from the top inward.  Comparable types
are replaced by their more specific type, and that refinement is substituted
through every occurrence of the same symbolic index.  Thus a late
\texttt{@} requirement can propagate through copies made earlier by
\texttt{SWAP} and \texttt{DUP}.  Incomparable boundary types make composition
undefined.  A sequence is composed from left to right, and the empty sequence
has the identity effect.  This operation extends earlier untyped stack-effect
composition \cite{poial1994}.

\subsection{Common branch contracts}

Alternative paths use the common-contract operation $p\meet q$.  It produces
one must-style interface rather than a set of path effects.  Its two basic
symbolic equations are
\[
 (a -- b)\meet(e\,a -- e\,b)=(e\,a -- e\,b),
\]
\[
 (e\,a -- e\,b)\meet(e\,c -- e\,d)
   =\bigl(e\,\min(a,c) -- e\,\max(b,d)\bigr).
\]
The first equation makes the implicit frame explicit: $(a--b)$ also acts as
$(e\,a--e\,b)$, so that is the common contract.  In the second equation,
$\min(a,c)$ means the more specific of two comparable input types, while
$\max(b,d)$ means the more general output type.  These are precision-order
operations, not numeric minimum and maximum.  Inputs narrow because a caller
must satisfy both paths; outputs widen because a client may assume only what
both paths guarantee.

Positions are aligned from the top.  An identity index survives only when both
paths preserve the same input-to-output relationship; a path-dependent
permutation is forgotten.  Unequal displayed depths are valid only when the
extra input and output depths are equal, because the omitted prefix must pass
through unchanged:
\[
 |L_p|-|L_q| = |R_p|-|R_q|.
\]
A different net stack change or an incomparable aligned type makes the meet
undefined (zero).

Write $\sem{P}$ for the inferred effect of a source segment $P$.  The selector
is composed separately for a conditional:
\[
 \sem{\texttt{IF }P\texttt{ ELSE }Q\texttt{ FI}}
   = (\texttt{flag}--)\seq(\sem{P}\meet\sem{Q}).
\]
Without \texttt{ELSE}, $Q$ is the identity effect.

\subsection{Finite repetition}
\label{sec:repeat-op}

For a repeated phrase with one-pass effect $p$, the prototype computes
\[
                  \pi_2(p)=p\meet(p\seq p),
\]
provided both the composition and common-contract operations are defined.
This one-versus-two comparison is finite and terminating.  It detects many
bodies that grow, shrink, or change the stack incompatibly on a second pass.

The notation has a narrow meaning.  It does not summarize every possible
iteration count.  One and two passes can agree after information has been
lost.  A later pass may then remove another correlation.  A full loop analysis
would propagate the loop-header state until that state stops changing.
$\pi_2$ does not perform that process.  It also proves neither termination nor
properties of concrete values.  It is only a local stack-stability test.

\section{Prototype Inputs}
\label{sec:control}

The prototype loads a type profile, a word-and-control profile, and the
example source.  Each example below introduces its primitive declarations at
the point of use rather than repeating them in a separate catalogue.

Control syntax is also profile-driven.  A conditional captures its two
branches and composes the selector with their common contract.  A
\texttt{BEGIN--WHILE--REPEAT} form captures a prefix and a body, while
\texttt{BEGIN--UNTIL} captures one body including the
final selector; the configured recipe applies $\pi_2$ to those segments.  The
profile therefore determines both how source is divided and which symbolic
operation combines each part.

With these profiles, the prototype accepts all twelve tests and the four
helper definitions.  Their inferred effects are derived once, in the following
sections.  The accepted top-level phrase is discussed with \texttt{test12}.

\section{Straight-Line Examples}
\label{sec:sequence-examples}

\subsection{Test1: stack shape without value evaluation}

\begin{quote}
\texttt{: test1  3 4 5 - + ;}
\end{quote}
Every literal has effect \effect{( -- n )}.  Naming fresh abstract results
$a$, $b$, and $c$, the stack evolves as
\[
 \epsilon
 \longrightarrow n_a
 \longrightarrow n_a\;n_b
 \longrightarrow n_a\;n_b\;n_c
 \longrightarrow n_a\;n_d
 \longrightarrow n_e.
\]
Subtraction consumes $n_b,n_c$ and produces a fresh $n_d$.  Addition consumes
$n_a,n_d$ and produces $n_e$.  No input is required.  One number is left:
\[
                      \texttt{test1}\;( -- n ).
\]
Concrete execution would calculate $4-5=-1$ and then $3+(-1)=2$.  The checker
does not perform that calculation.  Its domain contains the type \texttt{n},
but no integer constants or arithmetic expressions.  It establishes only
\effect{( -- n )}.

\subsection{Test2: a generic shuffle specialized to flags}

\begin{quote}
\texttt{: test2  OR FALSE SWAP ;}
\end{quote}
\texttt{OR} first requires two flags and leaves one fresh flag.  \texttt{FALSE}
adds another flag.  The generic inputs of \texttt{SWAP} then match these two
results.  Both inputs are refined from \texttt{x} to \texttt{flag}:
\[
 \texttt{flag}_a\;\texttt{flag}_b
 \xrightarrow{\texttt{OR}}
 \texttt{flag}_c
 \xrightarrow{\texttt{FALSE}}
 \texttt{flag}_c\;\texttt{flag}_d
 \xrightarrow{\texttt{SWAP}}
 \texttt{flag}_d\;\texttt{flag}_c.
\]
The inputs are consumed by \texttt{OR}; neither output is correlated with an
original input.  The reported interface therefore needs no visible indices:
\[
           \texttt{test2}\;(\texttt{flag flag -- flag flag}).
\]
Although one output is produced by the word named \texttt{FALSE}, the checker
has no constant-propagation domain and promises only its declared type.

\subsection{Test3: backward refinement through \texttt{SWAP} and
\texttt{DUP}}

\begin{quote}
\texttt{: test3  SWAP DUP @ ;}
\end{quote}
Let the initial suffix be $A\;B$, with $B$ topmost.  \texttt{SWAP} produces
$B\;A$ and \texttt{DUP} produces $B\;A\;A$.  The final \texttt{@} requires
the top copy of $A$ to be an \texttt{a-addr}.  Because both copies carry the
same symbolic index, the substitution also refines the untouched copy and the
original input:
\[
 A\;B
 \xrightarrow{\texttt{SWAP}} B\;A
 \xrightarrow{\texttt{DUP}} B\;A\;A
 \xrightarrow{\texttt{@}} B\;A\;X_{\mathrm{fresh}}.
\]
After normalized numbering, $A$ is \texttt{a-addr[2]} and $B$ is
\texttt{x[1]}.  Hence
\[
\begin{split}
 \texttt{test3}\;(&\texttt{a-addr[2] x[1]}\\
                  &\texttt{-- x[1] a-addr[2] x}).
\end{split}
\]
This is the key identity-flow example.  Tracking only stack depth would miss
the address requirement.  Tracking types without indices would also lose
information.  Such an analysis could refine only the consumed copy.  It could
not identify the surviving copy as the same aligned address.

\section{Branch Examples}
\label{sec:branch-examples}

\subsection{Test4: an explicit effect versus an implicit frame}

\begin{quote}
\texttt{: test4  IF 1 + FI ;}
\end{quote}
The true branch \texttt{1 +} pushes an \texttt{n}, then needs one additional
\texttt{n} below it for \texttt{+}.  Its inferred effect is
\effect{( n -- n )}.  The absent \texttt{ELSE} branch has the identity effect
\effect{( -- )}.  These displayed depths differ, but the empty branch preserves
an arbitrary deeper frame.  Making one frame cell explicit gives the
alignment
\[
\begin{array}{rcl}
 \text{true}  &:& (\texttt{n -- n}),\\
 \text{false} &:& (\underline{\texttt{n}} --
                       \underline{\texttt{n}}).
\end{array}
\]
Both paths preserve a one-cell numeric stack shape.  They do not preserve the
same value.  The false path returns the original number.  The true path returns
the result of addition.  No output identity is claimed.
After the leading \texttt{IF} consumes its selector, the complete definition is
\[
                    \texttt{test4}\;(\texttt{n flag -- n}).
\]
This example explains why comparing only the explicitly printed sides would be
too strict.  Stack effects denote transformations under an implicit frame.

\subsection{Test5: a common prefix survives, a branch-dependent
identity does not}

\begin{quote}
\texttt{: test5  IF DUP ELSE 0 FI ;}
\end{quote}
The true branch is
\effect{( x[1] -- x[1] x[1] )}.  The false branch explicitly has
\effect{( -- n )}, but it leaves the deeper frame unchanged.  Exposing one
frame cell aligns the branches as follows:
\[
\begin{array}{rcl}
 \text{true}  &:& (\texttt{x[1] -- x[1] x[1]}),\\
 \text{false} &:& (\texttt{x[1] -- x[1] n}).
\end{array}
\]
The lower output is the original input on both paths, so its index survives.
The top output is the same input on the true path but a fresh numeric literal
on the false path.  Its path-independent type is the more general type
\texttt{x}.  It cannot retain the input index.  The branch contract is thus
\[
                       (\texttt{x[1] -- x[1] x}).
\]
Composing the selector yields
\[
             \texttt{test5}\;(\texttt{x[1] flag -- x[1] x}).
\]
The result loses two facts.  First, \texttt{n} is widened to \texttt{x} at the
differing output.  Second, a correlation valid on only one path is discarded.

\subsection{Test6: compatible shape with no common permutation}

\begin{quote}
\texttt{: test6  IF ROT ELSE @ FI ;}
\end{quote}
Write the pre-branch suffix as $A\;B\;C$, with $C$ topmost.  The true branch
rotates it:
\[
                         A\;B\;C \longmapsto B\;C\;A.
\]
The false branch \texttt{@} displays only one input and one output.  Its
implicit two-cell frame makes the comparable transformation
\[
                         A\;B\;C \longmapsto A\;B\;D,
\]
where $C$ must be \texttt{a-addr} and $D$ is a fresh \texttt{x}.  At the input,
\texttt{ROT}'s generic requirement and \texttt{@}'s address requirement meet
at the more specific \texttt{a-addr}.  The input contract is consequently
\effect{( x x a-addr -- ... )}.

The output positions have the path pairs $(B,A)$, $(C,B)$, and $(A,D)$.  None
of these establishes one common source identity.  Their only common output
type is \texttt{x}, so the branch contract is
\[
                  (\texttt{x x a-addr -- x x x}).
\]
After selector consumption, the inferred definition is
\[
       \texttt{test6}\;(\texttt{x x a-addr flag -- x x x}).
\]
The checker accepts the equal net stack change.  It forgets which input reaches
which output.  An analysis that stores each path separately could retain both
permutations.  This checker stores one common contract.

\subsection{Test9: input and output variance in isolation}

\begin{quote}
\texttt{: test9  IF @ ELSE C@ FI ;}
\end{quote}
The two branch effects are
\[
 (\texttt{a-addr -- x})
 \qquad\text{and}\qquad
 (\texttt{c-addr -- char}).
\]
An input to the whole conditional must satisfy both address requirements.
Since \texttt{a-addr}$\preceq$\texttt{c-addr}, the common input is the more
specific \texttt{a-addr}.  An output client, however, needs a guarantee valid
for either \texttt{x} or \texttt{char}.  Since
\texttt{char}$\preceq$\texttt{x}, that guarantee is the more general
\texttt{x}.  Thus
\[
 (\texttt{a-addr -- x})\meet(\texttt{c-addr -- char})
       =(\texttt{a-addr -- x}),
\]
and selector composition gives
\[
               \texttt{test9}\;(\texttt{a-addr flag -- x}).
\]
This small example is the clearest demonstration of contract variance:
branch inputs are narrowed, whereas outputs are widened.

\subsection{Test10: merging effects inferred for helper words}

\begin{quote}
\texttt{: P INVERT HERE ! DP ! DP @ DP C@ SWAP ;}\\
\texttt{: Q SWAP ! DP @ DP ;}\\
\texttt{: test10 IF P ELSE Q FI ;}
\end{quote}
This example first checks two ordinary colon definitions and then uses their
inferred effects as branch contracts.  For \texttt{P}, let the initial suffix
be \texttt{x flag}.  \texttt{INVERT HERE !} consumes the new aligned address
and the flag result, leaving the lower \texttt{x}.  \texttt{DP !} then consumes
that cell and another aligned address.  Finally, \texttt{DP @} and
\texttt{DP C@} produce \texttt{x char}, which \texttt{SWAP} exchanges.  Thus
\[
                   \texttt{P}\;(\texttt{x flag -- char x}).
\]
Both outputs are fresh results of fetches; neither is an original input.

For \texttt{Q}, \texttt{SWAP !} consumes an \texttt{a-addr} and an arbitrary
cell.  \texttt{DP @ DP} then produces a fresh cell and a fresh aligned address:
\[
                   \texttt{Q}\;(\texttt{a-addr x -- x a-addr}).
\]
The common branch input narrows \texttt{x} against \texttt{a-addr} in the
lower position and \texttt{flag} against \texttt{x} in the upper position.
The two output pairs are \texttt{(char,x)} and \texttt{(x,a-addr)}; each
position therefore widens to \texttt{x}.  No identity is common to the paths:
\[
 \sem{\texttt{P}}\meet\sem{\texttt{Q}}
       = (\texttt{a-addr flag -- x x}).
\]
After \texttt{IF} consumes the selector, the complete result is
\[
       \texttt{test10}\;(\texttt{a-addr flag flag -- x x}).
\]
Unlike \texttt{test9}, this case shows that effects inferred for user-defined
words are immediately reusable and that variance is applied independently at
each aligned position.

\section{Repetition Examples}
\label{sec:loop-examples}

\subsection{Test7: one-versus-two exposes both flag requirements}

\begin{quote}
\texttt{: test7  BEGIN SWAP OVER WHILE INVERT REPEAT ;}
\end{quote}
The profile divides this structure into the repeated prefix
\texttt{SWAP OVER WHILE} and the repeated body \texttt{INVERT}.  Consider one
prefix pass on $A\;B$:
\[
 A\;B
 \xrightarrow{\texttt{SWAP}} B\;A
 \xrightarrow{\texttt{OVER}} B\;A\;B
 \xrightarrow{\texttt{WHILE}} B\;A.
\]
\texttt{WHILE} consumes the duplicate of the original top cell, so $B$ must be
a flag.  A second prefix pass transforms $B\;A$ back to $A\;B$ and tests $A$.
Consequently the two-pass effect requires both original cells to be flags.
The common summary preserves two flag-typed positions.  It cannot promise one
permutation.  One pass swaps the cells, while two passes restore them.  The
summary is therefore
\[
                        (\texttt{flag flag -- flag flag}).
\]
The second repeated region, \texttt{INVERT}, has the stable type effect
\effect{( flag -- flag )}.  Composing it with the prefix summary leaves the
same two-cell interface, giving
\[
          \texttt{test7}\;(\texttt{flag flag -- flag flag}).
\]
This is a useful rejection test for changing stack shape and incompatible
second-pass types.  It is not a proof that the reported identity-free summary
characterizes every number of loop iterations or that the loop terminates.

\subsection{Test8: a data-stack-neutral loop pass}

\begin{quote}
\texttt{: test8  BEGIN KEY 0= UNTIL ;}
\end{quote}
A complete repeated pass has the data-stack trace
\[
 \epsilon
 \xrightarrow{\texttt{KEY}} \texttt{char}
 \xrightarrow{\texttt{0=}} \texttt{flag}
 \xrightarrow{\texttt{UNTIL}} \epsilon.
\]
The loaded \texttt{x} contract for \texttt{0=} accepts \texttt{char}
directly.  The presentation's narrower \texttt{n} annotation would also accept
it because \texttt{char}$\preceq$\texttt{n}.  \texttt{UNTIL} consumes the
resulting flag.  One pass has the
identity data-stack effect.  Two passes have the same effect.  Their common
contract is
\[
                         \texttt{test8}\;(\texttt{--}).
\]
This conclusion concerns only the modeled data stack.  \texttt{KEY} still has
an I/O side effect.  The result does not promise a key that meets the exit
condition.  It also does not prove termination.

\subsection{Test11: a reusable body refines both loop cells}

\begin{quote}
\texttt{: BODY SWAP DUP @ 0= ;}\\
\texttt{: test11 BEGIN BODY UNTIL ;}
\end{quote}
The first three words of \texttt{BODY} repeat the identity flow from
\texttt{test3}.  The fetch result is a generic \texttt{x}; the loaded profile's
broad \effect{( x -- flag )} declaration for \texttt{0=} consumes it and
produces a selector.  The helper effect is
\[
 \texttt{BODY}\;(\texttt{a-addr[2] x[1] --
                         x[1] a-addr[2] flag}).
\]
\texttt{UNTIL} consumes the final flag.  One complete pass therefore swaps two
cells and requires the first one to be an aligned address:
\[
 p=(\texttt{a-addr[2] x[1] -- x[1] a-addr[2]}).
\]
On a second pass, the other original cell reaches \texttt{@} and is refined to
\texttt{a-addr}.  Hence $p\seq p$ requires two aligned addresses and restores
their order.  One pass swaps the origins while two passes preserve them, so no
input/output identity is common to both paths.  Their types are common:
\[
 p\meet(p\seq p)
   =(\texttt{a-addr a-addr -- a-addr a-addr}).
\]
This is the reported effect of \texttt{test11}.  It combines reusable inferred
identity information with the same two-pass refinement principle seen in
\texttt{test7}.

\subsection{Test12: a concrete boundary of one-versus-two}

\begin{quote}
\texttt{: BODY23 ROT INVERT DROP 0 ;}\\
\texttt{: test12 BEGIN BODY23 KEY? UNTIL ;}\\
\texttt{BODY23 BODY23 BODY23}
\end{quote}
Let the three inputs of \texttt{BODY23} be $A\;B\;C$, with $C$ topmost.
\texttt{ROT} produces $B\;C\;A$, so \texttt{INVERT} requires $A$ to be a
flag.  \texttt{DROP} removes the resulting flag, and the literal produces a
fresh \texttt{n}.  The reusable effect is
\[
 \texttt{BODY23}\;(\texttt{flag x[2] x[1] -- x[2] x[1] n}).
\]
In \texttt{test12}, \texttt{KEY?} produces the flag consumed by
\texttt{UNTIL}; those two words have no net data-stack effect.  Thus one
complete loop pass has the same effect as \texttt{BODY23}.  A second pass
requires $B$ to be a flag and maps $A\;B\;C$ to $C\;n\;n$.  Meeting one and two
passes retains the requirements that $A$ and $B$ are flags, widens the first
two output positions to \texttt{x}, and retains \texttt{n} at the top:
\[
 \texttt{test12}\;(\texttt{flag flag x -- x x n}).
\]
No index survives because the first output originates at $B$ after one pass
but at $C$ after two, while the remaining aligned outputs also lack a common
origin.

The final top-level phrase deliberately takes a third pass.  Its symbolic trace
is
\[
 \texttt{flag flag flag}
 \longrightarrow \texttt{flag flag n}
 \longrightarrow \texttt{flag n n}
 \longrightarrow \texttt{n n n}.
\]
The third pass introduces a requirement not present in the one-versus-two
summary: $C$ must also be a flag.  This does not contradict the rule, because
$\pi_2$ was defined as a bounded local comparison, not as an inductive
loop invariant.  It gives a concrete reason not to read the accepted
\texttt{test12} contract as valid for every iteration count.

\section{Discussion}

\subsection{What the examples establish}

The twelve tests isolate complementary properties of the symbolic domain.
\texttt{test1} shows that the checker infers shape and cell type without
constant evaluation.  \texttt{test2} shows subtype specialization of a generic
stack word.  \texttt{test3} demonstrates bidirectional information flow:
execution order is left to right, but a late requirement propagates backward
through shared symbolic identities.

The branch examples show why an effect is more than a pair of stack heights.
\texttt{test4} depends on the implicit frame.  \texttt{test5} retains only an
identity common to both paths.  \texttt{test6} loses a branch-dependent
permutation.  \texttt{test9} isolates input/output variance, and
\texttt{test10} applies that variance to effects inferred for helper words.
The loop examples expose both the value and the limit of the finite rule.  A
second pass refines another cell in \texttt{test7}; a balanced pass becomes
identity in \texttt{test8}; and \texttt{test11} carries identity information
through a reusable body.  Finally, \texttt{test12} and the following
three-\texttt{BODY23} phrase show directly why one-versus-two is not an
all-iteration invariant.

All reported outputs are inferred from primitive contracts.  The illustrative
comments in the presentation are not declarations.  This rule prevents a stale
comment from changing inference.  Semantic labels such as \texttt{value} or
``address contents'' are therefore outside the domain.  A richer profile would
be needed to represent them.

\subsection{Configuration as an analysis parameter}

The same source may receive a different verdict under another valid profile.
For example, one profile may classify \texttt{char} as numeric.  It accepts
the displayed \texttt{KEY 0=}.  Another profile may keep \texttt{char} and
\texttt{n} unrelated.  A numeric-only \texttt{0=} then rejects the phrase.
Even under the hierarchy used here, replacing the loaded broad
\effect{( x -- flag )} contract by \effect{( n -- flag )} would still accept
\texttt{test8} but reject \texttt{BODY}, because \texttt{@} guarantees only
\texttt{x}.  The target Forth system determines which policy is appropriate.
The profile makes that policy reviewable.

External control declarations provide similar flexibility.  They also enlarge
the trusted base.  A recipe might omit a selector effect.  It might also put
\texttt{UNTIL} outside its repeated segment.  Either error can produce a wrong
summary even when the three operations are correct.  The configuration should
therefore be reviewed with the checked code.

\section{Validity and Limitations}
\label{sec:limits}

Acceptance has a precise but bounded interpretation:
\begin{itemize}
  \item every checked linear boundary has enough modeled cells of comparable
        types;
  \item every checked branch admits one common contract under the defined
        alignment and variance rules;
  \item every checked repetition admits the finite $p\meet(p\seq p)$ summary;
  \item parser roles, delimiters, and compilation state fit the active profile.
\end{itemize}
It does not imply full Forth program correctness.  The principal threats to a
stronger interpretation are as follows.

\paragraph{Trusted declarations.}
Primitive and parser specifications are assumptions.  If \texttt{@},
\texttt{KEY}, or a control word is declared incorrectly, analysis of its
clients cannot make the result reliable.

\paragraph{Single effects, compaction, and comparable types.}
The dictionary stores one effect per word.  Branch merging also returns one
common contract.  It cannot retain a separate permutation for each path.
Normalization may discard an implicit one-to-one identity, so a stored derived
effect need not retain every correlation available in its construction trace.
Aligned types must be directly comparable.  The checker does not search for a
third common type.

\paragraph{Finite loop approximation.}
The one-versus-two comparison does not prove a result for every iteration
count.  A stronger analysis would maintain an abstract stack at the loop
header.  It would propagate the body until that stack stopped changing.
Another design would require a declared loop interface and check every back
edge against it \cite{cousot1977}.

\paragraph{Forth language coverage.}
Return-stack manipulation, unrestricted \texttt{EXECUTE}, recursion,
non-local \texttt{THROW}, native defining words, self-modifying dictionaries,
and unrestricted metaprogramming require additional abstractions.  The checker
should complement runtime tests, code review, and target-system knowledge.

The examples are explanatory cases, not a performance benchmark or a proof
that every accepted program is safe.  Each case makes one operation visible,
and the prototype's reported effects agree with \texttt{examples.txt.log}.

\section{Related Work}

Forth has long used stack-effect notation as an interface convention.  The
Forth standard defines data types and control semantics \cite{forth2012}.  It
does not make each source comment a static judgment.  Earlier work developed
algebraic composition and multiple effects for control flow
\cite{poial1994,poial1992}.  Later work added typed symbolic variables and
conservative control-flow operations \cite{poial2002,poial2006}.  The
prototype described here also loads parser and control behavior from a
profile.

Stoddart and Knaggs also studied type inference for stack languages
\cite{stoddart1991}.  Saabas and Uustalu gave compositional type systems for
stack-based low-level code \cite{saabas2006}.  Forth source can additionally
contain parser words and dictionary extensions.  The profile handles that
variability, but it becomes part of the trusted input.

Abstract interpretation offers a stronger treatment of loops
\cite{cousot1977}.  It propagates abstract states until they stop changing.
The one-versus-two rule performs only the first bounded comparison.  An
inferred loop effect must be read with that distinction in mind.

\section{Conclusion}
\leavevmode\par\vspace{-\baselineskip}
Typed symbolic stack effects turn a familiar Forth notation into executable,
configurable documentation.  Sequential composition checks touching stack
boundaries and propagates refined types through shared identities.  Branch
merging computes one conservative contract, narrowing inputs, widening outputs,
and preserving only path-independent correlations.  Repetition compares one
pass with two to obtain a useful but explicitly finite stability summary.

The tests provided make these mechanisms concrete.  Three initial linear
definitions show type and identity propagation.  Five conditionals show
implicit frames, input narrowing, output widening, information loss, and reuse
of inferred helper contracts.  Four loops show what the bounded repetition
rule accepts; the last is paired with a three-pass phrase that exposes exactly
what the rule does not establish.  The prototype derives every reported effect
under the selected profile.  It is a practical checker for the configured
subset of Forth.
It is not a proof of functional behavior or termination.  Stronger assurance
requires validated primitives, path-sensitive alternatives, less lossy
identity handling, and full loop analysis.

\begin{thebibliography}{99}
\footnotesize
\setlength{\itemsep}{0pt}
\setlength{\parsep}{0pt}
\raggedright

\bibitem{cousot1977}
P. Cousot and R. Cousot.
\newblock Abstract interpretation: A unified lattice model for static analysis
of programs by construction or approximation of fixpoints.
\newblock In \emph{Proceedings of POPL}, pp. 238--252, 1977.
\newblock doi:10.1145/512950.512973.
\newblock Public full text:
\url{https://www.di.ens.fr/~cousot/publications.www/CousotCousot-POPL-77-ACM-p238--252-1977.pdf}.

\bibitem{forth2012}
Forth 200x Standardisation Committee.
\newblock \emph{Forth 2012 Standard}.
\newblock Draft proposed standard, 10 November 2014.
\newblock Public full text:
\url{https://forth-standard.org/standard/intro}.

\bibitem{poial1994}
J. Pöial.
\newblock Algebraic specification of stack effects.
\newblock \emph{Forth Dimensions}, 16(4):18--20, 1994.
\newblock Public full text:
\url{https://www.forth.org/fd/FD-V16N4.pdf}.

\bibitem{poial1992}
J. Pöial.
\newblock Multiple stack-effects of Forth programs.
\newblock In \emph{1991 FORML Conference Proceedings, euroFORML'91},
pp. 400--406, 1992.
\newblock Public full text:
\url{https://www.complang.tuwien.ac.at/anton/euroforth/ef91/poial.pdf}.

\bibitem{poial2002}
J. Pöial.
\newblock Stack effect calculus with typed wildcards, polymorphism and
inheritance.
\newblock Presentation and proceedings abstract, 18th EuroForth Conference,
2002.
\newblock Public full text:
\url{https://www.complang.tuwien.ac.at/anton/euroforth/ef02/poial02.pdf}.

\bibitem{poial2006}
J. Pöial.
\newblock Typing tools for typeless stack languages.
\newblock In \emph{Proceedings of the 22nd EuroForth Conference},
pp. 40--46, 2006.
\newblock Public full text:
\url{https://www.complang.tuwien.ac.at/anton/euroforth/ef06/poial06.pdf}.

\bibitem{saabas2006}
A. Saabas and T. Uustalu.
\newblock Compositional type systems for stack-based low-level languages.
\newblock In \emph{Proceedings of CATS 2006}, CRPIT 51, pp. 27--39, 2006.
\newblock Public full text:
\url{https://web.archive.org/web/20120325100058id_/http://crpit.com/confpapers/CRPITV51Saabas.pdf}.

\bibitem{stoddart1991}
B. Stoddart and P. J. Knaggs.
\newblock Type inference in stack based languages.
\newblock In \emph{1991 FORML Conference Proceedings, euroFORML'91},
pp. 407--421, 1991.
\newblock Public full text:
\url{https://www.complang.tuwien.ac.at/anton/euroforth/ef91/stoddart.pdf}.

\end{thebibliography}

\end{document}

